% !TEX program = xelatex
\documentclass[11pt]{article}
\usepackage[a4paper,margin=24mm]{geometry}
\usepackage{amsmath,amssymb,amsthm,mathtools}
\usepackage{xcolor}
\usepackage[colorlinks=true,linkcolor=blue!50!black,citecolor=blue!50!black,urlcolor=blue!50!black]{hyperref}
\hypersetup{pdftitle={Boundary localization quotients and induced modules},pdfauthor={Discussion note}}
\newcommand{\C}{\mathbb C}
\newcommand{\Z}{\mathbb Z}
\newcommand{\g}{\widehat{\mathfrak{sl}}_2}
\newcommand{\ad}{\operatorname{ad}}
\newcommand{\im}{\operatorname{Im}}
\newcommand{\D}{\mathcal D}
\newcommand{\wt}{\operatorname{wt}}
\newcommand{\Span}{\operatorname{span}_{\C}}
\newtheorem{theorem}{Theorem}
\newtheorem{lemma}[theorem]{Lemma}
\newtheorem{corollary}[theorem]{Corollary}
\setlength{\emergencystretch}{2em}
\title{Boundary Localization Quotients and Induced Modules\\
  \large \(M_{0,0}(\psi)\cong T_{f_1}M_{-1,1}(\varphi)\)}
\author{Progress report}
\date{September 7, 2026}
\begin{document}
\maketitle

\section{Definitions, notation, and assumptions}

All vector spaces and tensor products are over \(\C\). We use the affine Lie
algebra without the degree derivation:
\[
 \g=(\mathfrak{sl}_2\otimes\C[t,t^{-1}])\oplus\C K,\qquad U=U(\g).
\]
Write \(x_i=x\otimes t^i\). For \(i,j\in\Z\), the brackets are
\begin{equation}\label{eq:brackets}
\begin{aligned}
 [e_i,f_j]&=h_{i+j}+i\delta_{i+j,0}K,&
 [h_i,h_j]&=2i\delta_{i+j,0}K,\\
 [h_i,e_j]&=2e_{i+j},&
 [h_i,f_j]&=-2f_{i+j},
\end{aligned}
\end{equation}
with \(K\) central and \([e_i,e_j]=[f_i,f_j]=0\);
\(\delta\) denotes the Kronecker delta.

\paragraph{Induced modules.}
Fix \(a,b\in\Z\) with \(a+b\ge0\), and put \(L=a+b+1\ge1\). Set
\[
 S_{a,b}=\Span\{K,\ h_i\ (i\ge0),\ e_i\ (i>a),\ f_j\ (j>b)\}.
\]
The brackets show that this is a Lie subalgebra, with
\[
 [S_{a,b},S_{a,b}]
 =\Span\{e_i\ (i>a),\ f_j\ (j>b),\ h_k\ (k\ge L+1)\}.
\]
Thus a character \(\chi:S_{a,b}\to\C\) is specified by arbitrary
\(c,\lambda_0,\ldots,\lambda_L\in\C\), through
\[
 \chi(K)=c,\quad \chi(h_i)=\lambda_i\ (0\le i\le L),\quad
 \chi(h_i)=0\ (i>L),\quad \chi(e_i)=\chi(f_j)=0,
\]
where the root vectors range over those in \(S_{a,b}\).
Set \(\lambda_i=0\) for \(i>L\). Define
\[
 \C_\chi=\C 1_\chi,\qquad
 M_{a,b}(\chi)=U\otimes_{U(S_{a,b})}\C_\chi,\qquad u_\chi=1\otimes1_\chi.
\]
Here \(s1_\chi=\chi(s)1_\chi\) for \(s\in S_{a,b}\).
In particular, \(h_0\in S_{a,b}\), so \(h_0u_\chi=\lambda_0u_\chi\).
These conventions agree with \cite[\S2.3]{FGXZ}.

\paragraph{PBW notation.}
Let \(\Gamma_r\) be the set of finite nondecreasing integer sequences
\(\alpha=(i_1,\ldots,i_k)\) with \(i_\nu\le r\), including the empty sequence.
Write
\[
 \ell(\alpha)=k,\qquad x(\alpha)=x_{i_1}\cdots x_{i_k},\qquad
 x(\varnothing)=1.
\]
Choose any total order on a basis of a vector-space complement of \(S_{a,b}\),
followed by a basis of \(S_{a,b}\). PBW gives a basis of the induced module
by the corresponding ordered complement monomials applied to \(u_\chi\).
The complement need not be a Lie subalgebra.

\paragraph{Localization and its quotient.}
For the boundary element \(F=f_b\), let \(\Sigma_F=\{F^r:r\ge0\}\).
The derivation \(\ad F\) is locally nilpotent on \(\g\), hence on \(U\).
The finite commutation formulas imply the left and right Ore conditions, so set
\[
 U_F=\Sigma_F^{-1}U,\qquad
 \D_FM=U_F\otimes_U M,\qquad T_FM=\D_FM/M.
\]
The quotient notation assumes that \(F\) acts injectively on \(M\), which
we verify below by PBW. Since \(M\) is a \(U\)-submodule of \(\D_FM\),
the quotient is a \(U\)-module.
The expression \(F^{-m}u+M\) always denotes a coset of a vector in the
localization; it does not assert that \(F^{-1}\) acts on the quotient.
A module \(V\) is smooth if each \(v\in V\) is annihilated by
\(e_r,h_r,f_r\) for all sufficiently large \(r\).

\paragraph{The basic case.}
Fix
\[
 M=M_{-1,1}(\varphi),\quad u=u_\varphi,\quad F=f_1,\quad
 T=T_FM,\quad w_m=F^{-m}u+M\quad(m\ge1),
\]
and write \(c=\varphi(K)\), \(\lambda_i=\varphi(h_i)\) for \(i\ge0\). Then
\begin{equation}\label{eq:cyclic}
\begin{gathered}
 e_i u=0\ (i\ge0),\qquad f_j u=0\ (j\ge2),\\
 h_i u=\lambda_i u\ (i\ge0),\qquad Ku=cu,\qquad \lambda_i=0\ (i\ge2).
\end{gathered}
\end{equation}
Define \(\psi:S_{0,0}\to\C\) by
\[
\begin{gathered}
 \psi(K)=c,\qquad \psi(h_0)=\lambda_0+2,\qquad
 \psi(h_1)=\lambda_1,\qquad \psi(h_i)=0\ (i\ge2),\\
 \psi(e_i)=\psi(f_i)=0\qquad(i\ge1),
\end{gathered}
\]
and put \(A=M_{0,0}(\psi)\), \(v=1\otimes1_\psi\).
The isomorphism requires \(\lambda_1\ne0\); the level \(c\) is arbitrary.
The additional condition \(c\ne-2\) is used only for simplicity.

For common PBW indices on the source and target, set
\[
 \mathcal X=\{X=e(\alpha)h(\beta)f(\gamma):
     \alpha,\beta\in\Gamma_{-1},\ \gamma\in\Gamma_0\},\qquad
 d(X)=2\ell(\alpha)+\ell(\beta).
\]
Write \(B_{m,X}=F^{-m}Xu+M\), and define vector subspaces
\begin{equation}\label{eq:filtration}
 T_{\le q}=\Span\{B_{m,X}:m\ge1,\ d(X)\le q\}\quad(q\ge0),
 \qquad T_{\le q}=0\quad(q<0).
\end{equation}
This is a vector-space filtration; its terms need not be \(U\)-submodules.

\section{The basic isomorphism}

\begin{theorem}\label{thm:basic}
If \(\lambda_1\ne0\), the assignment
\[
 \Phi:A=M_{0,0}(\psi)\longrightarrow T_{f_1}M_{-1,1}(\varphi),
 \qquad \Phi(xv)=xw_1\quad(x\in U)
\]
is the unique \(\g\)-module isomorphism sending \(v\) to \(w_1\).
This holds for every \(c\in\C\).
\end{theorem}

\subsection{Bases and commutation with negative powers}

Put \(F\) first in the complement PBW order for \(M\). Then
\[
 \{F^rXu:r\ge0,\ X\in\mathcal X\}
\]
is a basis of \(M\). Consequently \(M\) is free over \(\C[F]\), and \(F\)
acts injectively.
The localization has basis \(\{F^rXu:r\in\Z,\ X\in\mathcal X\}\):
common denominators give spanning, and multiplying a finite relation on
the left by a sufficiently large power of \(F\) reduces independence to PBW in \(M\).
It follows that
\begin{equation}\label{eq:bases}
\begin{aligned}
 &\{B_{m,X}:m\ge1,\ X\in\mathcal X\}&&\text{is a basis of }T,\\
 &\{Xe_0^pv:p\ge0,\ X\in\mathcal X\}&&\text{is a basis of }A.
\end{aligned}
\end{equation}
For the second basis, place \(e_0\) after all factors of \(X\) and before
the inducing subalgebra. In particular, \(w_1\ne0\).

Put \(D=\ad F\). The brackets give
\begin{equation}\label{eq:ad}
\begin{aligned}
 D(f_n)&=0,& D(h_n)&=2f_{n+1},\\
 D(e_n)&=-(h_{n+1}+n\delta_{n+1,0}K),&
 D^2(e_n)&=-2f_{n+2},\qquad D^3(e_n)=0.
\end{aligned}
\end{equation}
For \(x\in U\), \(m\ge1\), the following finite sum is an identity in \(U_F\):
\begin{equation}\label{eq:commute}
 xF^{-m}=\sum_{j\ge0}\binom{m+j-1}{j}F^{-m-j}D^j(x).
\end{equation}
For \(m=1\), multiply the right-hand side on the right by \(F\) and use
\(D^j(x)F=FD^j(x)-D^{j+1}(x)\); the sum telescopes to \(x\).
Iteration gives the formula for general \(m\), using
\(\sum_{r=0}^j\binom{m+r-1}{r}=\binom{m+j}{j}\).

\subsection{Well-definedness and the module homomorphism}

Equation \eqref{eq:commute} gives
\begin{equation}\label{eq:actions}
\begin{aligned}
 f_nF^{-m}&=F^{-m}f_n,\\
 h_nF^{-m}&=F^{-m}h_n+2mF^{-m-1}f_{n+1},\\
 e_nF^{-m}&=F^{-m}e_n
  -mF^{-m-1}(h_{n+1}+n\delta_{n+1,0}K)
  -m(m+1)F^{-m-2}f_{n+2}.
\end{aligned}
\end{equation}
Taking \(m=1\) and applying \eqref{eq:cyclic} yields
\[
\begin{gathered}
 e_nw_1=f_nw_1=0\ (n\ge1),\qquad Kw_1=cw_1,\\
 h_0w_1=(\lambda_0+2)w_1,\qquad h_1w_1=\lambda_1w_1,\qquad
 h_nw_1=0\ (n\ge2).
\end{gathered}
\]
Indeed, \(f_1w_1=u+M=0\). For \(n\ge1\), the correction terms for \(e_n\)
vanish because \(h_{n+1}u=f_{n+2}u=0\); those for \(h_n\) vanish because
\(f_{n+1}u=0\). The correction for \(h_0\) is
\(2F^{-2}f_1u+M=2w_1\).
Thus \(sw_1=\psi(s)w_1\) for all \(s\in S_{0,0}\).

Let \(\psi_U:U(S_{0,0})\to\C\) be the unital algebra character extending
\(\psi\). Applying the preceding relation to finite products and linear
combinations gives \(aw_1=\psi_U(a)w_1\) for every \(a\in U(S_{0,0})\).
The linear map
\[
 U\otimes_\C\C_\psi\longrightarrow T,\qquad
 x\otimes z1_\psi\longmapsto z\,xw_1
\]
kills every balancing relation, since
\[
 xa\otimes z1_\psi-x\otimes a(z1_\psi)
 \longmapsto z\,x(aw_1-\psi_U(a)w_1)=0.
\]
It therefore descends to a well-defined \(\Phi:A\to T\).
For \(y,x\in U\),
\(\Phi(y(x\otimes1_\psi))=yxw_1=y\Phi(x\otimes1_\psi)\),
so \(\Phi\) is \(U\)-linear. Uniqueness follows from \(A=Uv\).
No nonvanishing assumption on \(\lambda_1\) has been used at this stage.

\subsection{Explicit preimages of all pure negative powers}

Taking \(n=0\) in \eqref{eq:actions}, and using
\(e_0u=f_2u=0\), \(h_1u=\lambda_1u\), gives
\begin{equation}\label{eq:tower}
 e_0w_m=-m\lambda_1w_{m+1},\qquad
 e_0^pw_1=(-1)^pp!\lambda_1^pw_{p+1}\quad(p\ge0).
\end{equation}
Since \(\lambda_1\ne0\), we have
\begin{equation}\label{eq:preimage}
 w_m=\Phi(q_m),\qquad
 q_m=\frac{(-1)^{m-1}}{(m-1)!\lambda_1^{m-1}}e_0^{m-1}v\quad(m\ge1).
\end{equation}
This explicitly gives a preimage of each \(f_1^{-m}u+M\), without assuming
surjectivity.

\subsection{A strictly decreasing integer}

\begin{lemma}\label{lem:lower}
For \(m\ge1\) and \(X\in\mathcal X\),
\begin{equation}\label{eq:lower}
 Xw_m=B_{m,X}+R_{m,X},\qquad R_{m,X}\in T_{\le d(X)-1},
\end{equation}
where \(R_{m,X}\) is a finite linear combination.
\end{lemma}

\begin{proof}
Assign nonnegative weights
\[
 \wt(e_i)=2,\qquad \wt(h_i)=1,\qquad \wt(f_i)=0,\qquad \wt(K)=0,
\]
and filter \(U\) by total word weight, writing \(U_{\le q}=0\) for \(q<0\).
Each bracket has weight at most the sum of its input weights:
for \([h,e]\), \([e,f]\), \([h,f]\), and \([h,h]\), the output weights
are at most \(2,1,0,0\), whereas the input sums are \(3,2,1,2\).
Thus PBW reordering never increases weight. Neither does evaluating
the inducing factors on \(u\), where they become scalars or zero.

Equation \eqref{eq:ad} and the Leibniz rule give
\[
 D(U_{\le q})\subset U_{\le q-1},\qquad
 D^j(X)\in U_{\le d(X)-j}.
\]
In particular \(D^j(X)=0\) for \(j>d(X)\). Expanding in the PBW basis of \(M\),
we obtain a finite sum
\[
 D^j(X)u=\sum_{r\ge0,\,X'\in\mathcal X}c_{j,r,X'}F^rX'u,\qquad
 c_{j,r,X'}\ne0\ \Longrightarrow\ d(X')\le d(X)-j.
\]
Substitute this into \eqref{eq:commute}. The \(j=0\) term is exactly \(B_{m,X}\).
For \(j\ge1\), a term with \(r\ge m+j\) belongs to \(M\) and vanishes in \(T\).
Every surviving term is \(B_{m+j-r,X'}\), with \(d(X')<d(X)\).
This proves \eqref{eq:lower}.
\end{proof}

\subsection{Surjectivity by strong induction}

For \(d\ge0\), prove the statement
\[
 P(d):\quad B_{m,X}\in\im\Phi
 \quad\text{for every }m\ge1\text{ and every }X\in\mathcal X
 \text{ with }d(X)\le d.
\]
The induction hypothesis includes all denominator exponents simultaneously;
an increase in the denominator exponent of a correction term is therefore harmless.

For \(d=0\), we have \(X=f(\gamma)\), which commutes with \(F\). Hence
\[
 B_{m,X}=Xw_m=\Phi(Xq_m).
\]
The base case includes every \(f\)-monomial, not just \(X=1\).

Suppose \(d\ge1\) and \(P(d-1)\) holds. Fix \(m\ge1\) and \(d(X)=d\).
Lemma~\ref{lem:lower} gives a finite expansion
\[
 \Phi(Xq_m)=B_{m,X}+\sum_{\nu=1}^{s}c_\nu B_{m_\nu,X_\nu},
 \qquad m_\nu\ge1,\quad d(X_\nu)\le d-1.
\]
By induction, choose \(a_\nu\in A\) with
\(\Phi(a_\nu)=B_{m_\nu,X_\nu}\). Then
\[
 B_{m,X}=\Phi\left(Xq_m-\sum_{\nu=1}^{s}c_\nu a_\nu\right).
\]
This proves \(P(d)\) and recursively constructs a preimage.
Each recursion involves finitely many terms and strictly decreases the
nonnegative integer \(d\), so it terminates.
The target basis in \eqref{eq:bases} proves surjectivity.

For example, \(d(h_{-1})=1\), and
\[
 h_{-1}w_m=B_{m,h_{-1}}+2mB_{m+1,f_0}.
\]
The correction has degree zero, giving the concrete preimage
\[
 B_{m,h_{-1}}=\Phi(h_{-1}q_m-2mf_0q_{m+1}).
\]
This also illustrates why the induction must include all \(m\).

\subsection{Injectivity by the highest degree}

By \eqref{eq:tower} and Lemma~\ref{lem:lower},
\begin{equation}\label{eq:leading}
 \Phi(Xe_0^pv)\equiv(-1)^pp!\lambda_1^p B_{p+1,X}
       \pmod{T_{\le d(X)-1}}.
\end{equation}
Let \(0\ne a=\sum_{X,p}a_{X,p}Xe_0^pv\in A\), a finite sum, and let
\(d=\max\{d(X):a_{X,p}\ne0\}\). In \(T_{\le d}/T_{\le d-1}\),
\[
 \Phi(a)+T_{\le d-1}
 =\sum_{\substack{X,p\\d(X)=d}}
    a_{X,p}(-1)^pp!\lambda_1^p
       (B_{p+1,X}+T_{\le d-1}).
\]
These cosets are linearly independent by \eqref{eq:bases}, and
\((-1)^pp!\lambda_1^p\ne0\). Thus \(\Phi(a)\ne0\), proving injectivity
and Theorem~\ref{thm:basic}.
Only finite linear combinations are used; the filtration terms need not be
finite-dimensional.

\section{Direct extensions}

\subsection{The boundary of \texorpdfstring{\(M_{-1,N}\)}{M(-1,N)}}

\begin{theorem}\label{thm:N}
Let \(N\ge1\), \(M=M_{-1,N}(\varphi)\), \(u=u_\varphi\),
\(c=\varphi(K)\), and \(\lambda_i=\varphi(h_i)\).
Assume \(\lambda_N\ne0\). Define \(\psi:S_{0,N-1}\to\C\) by
\[
\begin{gathered}
 \psi(K)=c,\qquad \psi(h_i)=\lambda_i+2\delta_{i,0}\quad(i\ge0),\\
 \psi(e_i)=0\quad(i\ge1),\qquad \psi(f_j)=0\quad(j\ge N).
\end{gathered}
\]
Then there is an isomorphism
\[
 \Phi_N:M_{0,N-1}(\psi)\xrightarrow{\ \sim\ }T_{f_N}M_{-1,N}(\varphi),
 \qquad v\longmapsto f_N^{-1}u+M.
\]
\end{theorem}

\begin{proof}
Put \(F=f_N\), \(w_m=F^{-m}u+M\). In the commutation formulas, replace
\(n+1,n+2\) by \(n+N,n+2N\), respectively. Now
\(e_iu=0\) for \(i\ge0\), \(f_ju=0\) for \(j>N\), and
\(\lambda_i=0\) for \(i>N\).
Consequently \(e_iw_1=0\) for \(i\ge1\) and \(f_jw_1=0\) for \(j\ge N\);
at the boundary \(f_Nw_1=u+M=0\). Also
\[
 h_iw_1=(\lambda_i+2\delta_{i,0})w_1\quad(i\ge0),\qquad Kw_1=cw_1.
\]
These relations give the well-defined homomorphism \(\Phi_N\) by the
balancing argument. Since \(e_0u=f_{2N}u=0\),
\[
 e_0w_m=-m\lambda_Nw_{m+1},\qquad
 w_m=\Phi_N\left(\frac{(-1)^{m-1}e_0^{m-1}v}
                         {(m-1)!\lambda_N^{m-1}}\right).
\]
Replace \(\mathcal X\) by
\[
 \mathcal X_N=\{e(\alpha)h(\beta)f(\gamma):
      \alpha,\beta\in\Gamma_{-1},\ \gamma\in\Gamma_{N-1}\}.
\]
The target and source bases are \(\{F^{-m}Xu+M\}\) and \(\{Xe_0^pv\}\).
The same degree \(d(X)=2\ell(\alpha)+\ell(\beta)\) is strictly lowered by
\(\ad f_N\), while reordering and inducing relations do not increase it.
Thus Lemma~\ref{lem:lower} applies unchanged. The same strong induction
proves surjectivity, and the highest-degree coefficients
\((-1)^pp!\lambda_N^p\ne0\) prove injectivity.
\end{proof}

\subsection{A single step at an arbitrary boundary}

\begin{corollary}\label{cor:general}
Let \(a,b\in\Z\), \(a+b\ge0\), \(L=a+b+1\), \(M=M_{a,b}(\chi)\),
and \(u=u_\chi\). If \(\lambda_L=\chi(h_L)\ne0\), then
\[
 M_{a+1,b-1}(\chi^+)\xrightarrow{\ \sim\ }T_{f_b}M,\qquad
 v\longmapsto f_b^{-1}u+M,
\]
where \(\chi^+\) vanishes on all root vectors in \(S_{a+1,b-1}\), and
\[
 \chi^+(K)=c,\qquad \chi^+(h_i)=\lambda_i+2\delta_{i,0}\quad(i\ge0).
\]
\end{corollary}

\begin{proof}
Set \(F=f_b\), \(Y=e_{a+1}\), \(w_m=F^{-m}u+M\).
For \(i\ge a+2\), commuting \(e_i\) past \(F^{-1}\) gives a Cartan index
\(i+b\ge L+1\), so the Cartan and central terms vanish on \(u\).
The remaining correction vanishes because \(i+2b>b\), hence
\(f_{i+2b}u=0\).
The other inducing relations are
\[
 f_jw_1=0\ (j\ge b),\quad
 h_iw_1=(\lambda_i+2\delta_{i,0})w_1\ (i\ge0),\quad Kw_1=cw_1.
\]
Thus the homomorphism is well defined. Since \(Yu=0\), \(L>0\), and
\(f_{a+1+2b}u=0\), we obtain \(Yw_m=-m\lambda_Lw_{m+1}\).
Use
\[
 \mathcal X_{a,b}=\{e(\alpha)h(\beta)f(\gamma):
    \alpha\in\Gamma_a,\ \beta\in\Gamma_{-1},\ \gamma\in\Gamma_{b-1}\}
\]
and the source basis \(\{XY^pv\}\).
The derivation \(\ad f_b\) strictly lowers
\(d(X)=2\ell(\alpha)+\ell(\beta)\), so the same surjectivity induction and
highest-degree comparison apply, with coefficient
\((-1)^pp!\lambda_L^p\ne0\).
\end{proof}

\subsection{Smoothness, simplicity, and the assumptions}

The induced modules above are smooth. Indeed, move sufficiently high modes
past a fixed finite PBW word; all resulting noncentral modes still have
sufficiently high indices and annihilate the inducing vector.
The bound depends on the word. The isomorphisms therefore also establish
smoothness of the localization quotients.

By \cite[Theorem~1.1]{FGXZ}, for \(a+b\ge0\),
\[
 M_{a,b}(\chi)\text{ is simple}\quad\Longleftrightarrow\quad
 \lambda_{a+b+1}\ne0\text{ and }c\ne-2.
\]
Under the hypothesis \(\lambda_L\ne0\) of Corollary~\ref{cor:general},
the boundary quotient is simple if and only if \(c\ne-2\).
The isomorphism remains valid at the critical level \(c=-2\);
its injectivity was proved without using simplicity.

If \(\lambda_1=0\), the basic homomorphism still exists, but
\(\Phi(e_0v)=e_0w_1=0\), whereas \(e_0v\ne0\) by the source PBW basis.
Thus the specified map is not an isomorphism.

\begin{thebibliography}{9}
\bibitem{FGXZ}
V.~Futorny, X.~Guo, Y.~Xue and K.~Zhao,
\emph{Smooth representations of affine Kac--Moody algebras},
arXiv:2404.03855v2 (2024), \S2.3 and Theorem~1.1.
\url{https://arxiv.org/abs/2404.03855v2}.
\end{thebibliography}
\end{document}
